% Source: https://artofproblemsolving.com/community/c3h3242498_the_polymath_competition
% Stable identity: aops-topic-3242498-q30
Equivalently, whenever $N\mid f(X)$, we require $N\mid X$. For even $X$, this always holds because $X=2f(X)$. If $3\mid N$, then for odd $X$ the number $3X+1$ is not divisible by $3$, hence not by $N$. Thus every multiple of $3$ works. Also $N=1$ works vacuously.
Now take $N>1$ coprime to $3$. There is a residue $r$ with $3r\equiv-1\pmod N$, and it is nonzero modulo $N$. We can choose a positive odd $X\equiv r\pmod N$: if $N$ is odd, adding $N$ changes parity; if $N$ is even, the congruence already forces $r$ odd. Then $N\nmid X$ but $N\mid3X+1=f(X)$, a contradiction. Therefore exactly $\boxed{N=1\text{ or }3\mid N}$ work.
